\documentclass[11pt]{ltjsarticle}
\usepackage{luatexja-fontspec}
\usepackage{amsmath,amssymb}
\usepackage[margin=22mm]{geometry}
\usepackage{graphicx}
\usepackage{longtable,booktabs,array}
\usepackage{calc}
\usepackage{xcolor}
\usepackage{hyperref}
\hypersetup{colorlinks=true,linkcolor=blue,urlcolor=blue,citecolor=blue}
\makeatletter
\def\maxwidth{\ifdim\Gin@nat@width>\linewidth\linewidth\else\Gin@nat@width\fi}
\makeatother
\setkeys{Gin}{width=\maxwidth,keepaspectratio}
\setcounter{secnumdepth}{0}
\setlength{\parindent}{0pt}
\setlength{\parskip}{0.5em}
\providecommand{\tightlist}{\setlength{\itemsep}{0pt}\setlength{\parskip}{0pt}}
\providecommand{\pandocbounded}[1]{#1}
\providecommand{\real}[1]{#1}
\newcounter{none}
\begin{document}
\section{On the Logarithmic Map of Scale and the Separation of Scale and
Quantization in a High-Symmetry
Limit}\label{on-the-logarithmic-map-of-scale-and-the-separation-of-scale-and-quantization-in-a-high-symmetry-limit}

\subsection{--- A Limited Verification Report from a Minimal Axiom
System
(Revised)}\label{a-limited-verification-report-from-a-minimal-axiom-system-revised}

\textbf{Author:} Noriaki Kihara (ORCID: 0009-0004-6753-4020)
\textbf{Date:} 18 June 2026 (revised) \textbf{DOI (this version, v2):}
10.5281/zenodo.20742386 (v1: 10.5281/zenodo.20740842) \textbf{Concept
DOI (all versions):} 10.5281/zenodo.20740841 \textbf{Zenodo:}
https://zenodo.org/records/20742386 \textbf{v2 revision (2026-06-18):}
wording fix in §6.3 (what is preserved is the product of observables,
not the product of fluctuations); added a list of related supplements.

\begin{center}\rule{0.5\linewidth}{0.5pt}\end{center}

\subsection{Note to the reader (please read
first)}\label{note-to-the-reader-please-read-first}

This report does not claim the ``discovery'' of a new general fact.
Because it starts from a deliberately limited set of premises, it is
easily misread on a first pass. Most misreadings arise from
\textbf{blurring the system-interior (the unknowable, anonymous internal
viewpoint) and the system-exterior (the observer's viewpoint)}.
Moreover, in this system the question ``which side is real?'' is
ill-posed (see the note in §0.5); the objective anchors are placed not
on ``reality'' but on fixed points and conserved quantities. To prevent
misreading, each statement is, in principle, prefixed with a tag
indicating whether it is an interior or an exterior statement. Please
read §0.5 (interior and exterior) and §1 (conventions) first.

\begin{center}\rule{0.5\linewidth}{0.5pt}\end{center}

\subsection{Abstract}\label{abstract}

This report begins with a naive question about the conjugate relation of
two quantities, νλ=k (a quantity read off by an observer). This rigid
relation alone generates no structure. As the only crevice through which
structure can enter, we posit that the integer-valued ν carries an
intrinsic fluctuation, and that the system possesses high anonymity and
high symmetry. In a limited regime requiring extremely high symmetry and
anonymity, we report a \textbf{result within a limited scope}: a squared
quantity representing scale can be mapped, through a logarithmic
connection, onto a flat additive structure (a flat hyperplane).

Introducing the fundamental quantity as the logarithm of a squared
quantity, N=log\textsubscript{2}ν, the conjugate relation νλ=k becomes an additive
identity N\_ν+N\_λ=K, and the sum of squares Σν²=R² (the g side)
representing scale is mapped, in the self-dual fixed point (k=1) and the
fully symmetric limit, onto a flat linear sum ΣN=K. What matters is
that, while this logarithmic map flattens scale (the sum of squares, g)
and removes it as a gauge, the fluctuation Δ (the product / area / ω)
remains invariant on that flat surface. That is, the map
\textbf{separates scale from quantization}: scale resides in g and is
flattenable, whereas quantization resides in ω and is invariant under
the map. Here R is the scale (magnitude) of the system, and is neither a
curvature tensor nor the radius of a sphere.

This report treats the magnitude of the fluctuation, ±1/2, as a
\textbf{posit, not a derivation} (reasons in §6). Whether the result
remains a special case of known frameworks (Weyl transformation,
conformal flatness, etc.) cannot be determined here, and an exhaustive
search for counterexamples has not been carried out. The purpose of this
report is confined to recording this limited structure.

\begin{center}\rule{0.5\linewidth}{0.5pt}\end{center}

\subsection{§0 Motivation (author's
note)}\label{motivation-authors-note}

The starting point of this report is a naive question about the
conjugate relation νλ=k.

{[}Exterior (observer){]} Writing νλ=k gives λ=k/ν: as ν grows, λ
shrinks. This is a rigid see-saw relation read off from the outside.

{[}Exterior (observer){]} But this rigid relation alone yields no
structure from the two quantities ν and λ. There is merely a single
proportionality.

Where, then, can structure enter? This report loosens two points.

{[}Discrete structure{]} First, ν is discrete (integer-valued) and has
no continuous fractional part. It is not that ``a continuous value
really exists but appears integral due to the limits of observational
precision''; rather, \textbf{the fractional part does not exist} (this
``does not exist'' is a statement about the system-exterior, i.e., the
ontic level, and is distinct from the interior indistinguishability
discussed later). Yet the integer value is not fixed: it carries an
intrinsic fluctuation. This is not observational error but a fluctuation
intrinsic to the system. (Integrality itself is not a claim about
reality, but a consequence of discreteness; its derivation is in §3.)

{[}Interior{]} Second, in νλ=k there is no name distinguishing ν from λ.
From within the system, which is which cannot be told. The system thus
possesses high anonymity.

This report deliberately builds the system under this high anonymity and
high symmetry. Ordinarily, imposing such strong anonymity and symmetry
constrains the degrees of freedom and makes structure hard to arise. The
aim of this report is to see what nonetheless stands up.

{[}Exterior (observer) / Interior{]} One last point. νλ=k is, after all,
a quantity read off from the observer's viewpoint, exterior to the
system. Apart from it, might there be an invariant (a conserved
quantity) or a fixed point (a stable point) intrinsic to the interior?
This is the question running through the report.

\begin{center}\rule{0.5\linewidth}{0.5pt}\end{center}

\subsection{§0.5 Author's note: interior and exterior (the unknowable
and
invariants)}\label{authors-note-interior-and-exterior-the-unknowable-and-invariants}

Throughout, this report strictly distinguishes the system-interior from
the system-exterior. Blurring this causes most of the statements below
to be misread.

\begin{itemize}
\tightlist
\item
  {[}Interior{]} The interior is the viewpoint of an anonymous element
  (one with no name, color, or position). What cannot be distinguished
  within the interior is treated as equivalent to \emph{null} (nothing)
  within the interior. From the interior, exterior quantities cannot, in
  principle, be observed directly (except when an exterior quantity
  manifests inside as an effect; see Note 1).
\item
  {[}Exterior{]} The exterior is the viewpoint of one who observes from
  outside. The observer obtains only a limited set of observable
  quantities, and does not necessarily reach the interior invariants
  (fixed points, conserved quantities). Hence the observer cannot assert
  what is (or is not) inside the system.
\end{itemize}

\textbf{{[}Important note: reality cannot be asked.{]}} This system
rests on a dual / conjugate relation (νλ=k). Hence the question ``which
of ν and λ is real?'' is ill-posed --- declaring one real makes the
reality of the other ambiguous. Nor do we deny interior reality. What
carries meaning is not ``which is real'' but what is held invariant
under the duality --- the \textbf{fixed points (k=1, ν=1) and conserved
quantities (the product δ\textsubscript{1}δ\textsubscript{2}=1/2, the sum of squares Σν²)}. In what
follows, all objective anchors are stated as fixed points / conserved
quantities, and no claim is made about the locus of reality.

Under this distinction, we derive the most basic counting from the
interior viewpoint. The only thing registrable from the interior is
``whether there is something other than oneself, or not.''

\begin{itemize}
\tightlist
\item
  {[}Interior{]} The case 0: there is no subject to judge a state. Hence
  0 cannot be distinguished from null. This only states interior
  indistinguishability and does not exclude 0 as an integer-structure
  entity.
\item
  {[}Interior{]} The case 1: the sole element has no ``other than
  itself.'' Being anonymous, ``I exist'' cannot even stand without a
  contrast against an other. Hence, in the interior, 1 cannot be
  distinguished from 0. This does not exclude 1 as an integer-structure
  entity.
\item
  {[}Interior{]} The case 2: for the first time there is an ``other than
  oneself.'' Through it, ``I exist'' stands. Two is the minimum at which
  anything can be known from the interior.
\item
  {[}Interior{]} The case 2-or-more: the other carries no color, name,
  or position, and so cannot be counted. Hence, in the interior, 2 and
  2-or-more cannot be distinguished. This does not exclude 2-or-more as
  an integer-structure entity.
\end{itemize}

{[}Interior{]} Therefore the interior viewpoint collapses to a binary
--- ``no other'' (0 and 1 indistinguishable) and ``an other present'' (2
and 2-or-more indistinguishable).

{[}Conserved quantity{]} On the other hand, as a structure of conserved
quantities / fixed points, the integer sequence 0,1,2,3,\ldots{} stands.
The interior collapse to a binary is a collapse by unknowability, and
does not negate this integer structure.

{[}Interior / Exterior{]} In one line: what is unknowable in the
interior is, in the interior, equivalent to null. But it cannot be
negated from the exterior either --- what is unknowable is not even an
object of negation.

\textbf{Note 1 (an example of an exterior quantity manifesting inside as
an effect):} spatial curvature in a multidimensional spacetime is such
an example. Curvature itself is an exterior quantity, and the interior
cannot perceive its direction; but its effect --- for instance, the sum
of the interior angles of a triangle deviating from 180° --- can be
perceived from the interior. In the present system, however, this
example does not yet clearly arise, and curvature lies outside the scope
of this report (§7, §8).

\begin{center}\rule{0.5\linewidth}{0.5pt}\end{center}

\subsection{§1 Conventions (how to read
interior/exterior)}\label{conventions-how-to-read-interiorexterior}

We now enter the formal description. For readability, two conventions
are set.

\begin{enumerate}
\def\labelenumi{\arabic{enumi}.}
\tightlist
\item
  \textbf{{[}Tag{]}} Each statement is, in principle, prefixed with a
  tag indicating whether it belongs to the interior (the unknowable,
  anonymous internal viewpoint) or the exterior (the observer's
  viewpoint). Anchors held invariant under the duality are marked as
  {[}Fixed point{]}, {[}Conserved quantity{]}, or {[}Discrete
  structure{]}; no claim is made about the locus of reality.
\item
  \textbf{{[}Verb discipline{]}} When stating interior
  indistinguishability, we do not write ``X does not exist'' but ``X
  cannot be distinguished in the interior.'' ``Does not exist'' is
  allowed only for objective statements about discrete structure / fixed
  points / conserved quantities, and is not used for interior
  indistinguishability.
\end{enumerate}

\textbf{{[}Declaration{]}} Among the axioms and derivations below, all
statements about 0, 1, 2 are statements about interior
distinguishability, and do not negate the integer structure
(discreteness / fixed points / conserved quantities).

\begin{center}\rule{0.5\linewidth}{0.5pt}\end{center}

\subsection{§2 Axioms (assumptions)}\label{axioms-assumptions}

The system rests on the following axioms. No external premises or
references to prior literature are used. Physical interpretation is kept
separate; we restrict ourselves to statements about structure.

\textbf{Axiom 1 (existence of difference):} There is difference. Only
``there is another quantity'' can be said. What differs, which differs
from which, and how many there are cannot be said (anonymity). Speaking
of difference itself presupposes difference, so this is
self-referential.

\textbf{Axiom 2 (the limit of anonymity = 2):} {[}Interior{]} The
minimum at which difference stands is two. By anonymity, the maximum is
also two. In the interior it closes on 0, 1, 2, with 2 as the ceiling
(by the derivation in §0.5: 0 and 1 cannot be distinguished, and 2 and
2-or-more cannot be distinguished). Treating 3-or-more \emph{as}
3-or-more is counting = labeling = the genesis of dimension, and is a
breaking of anonymity. This is the limit of interior distinguishability,
and does not exclude 3-or-more as an integer-structure entity.

\textbf{Axiom 3 (fluctuation):} Each quantity has an intrinsic
fluctuation δ. δ is an attribute of that quantity and is a scalar. The
magnitude ±1/2 of the fluctuation is a posit, not a derivation (§6).

\textbf{Axiom 4 (conservation of the product: the edge of being /
non-being):} {[}Interior{]} The product of the two fluctuations is
conserved: δ\textsubscript{1}δ\textsubscript{2}=1/2. This is the invariant of the edge of the binary of
existence (``being / non-being'') (product, hyperbolic SO(1,1), area),
and corresponds to the imaginary part ω of the complex inner product
introduced later. (Which one is taken as the imaginary part is, by the
principle of anonymity, itself a matter of choice and is arbitrary.)

\textbf{Axiom 5 (conservation of the sum of squares: the binding of
difference):} {[}Interior{]} The system is bounded and symmetric
(anonymity invariant under relabeling). This makes an invariant binding
the two on equal footing stand: Σν²=R² (sum of squares, SO(2) symmetry,
scale R = distance g). Here SO(2) refers only to the symmetry preserving
the sum of squares, and does not assume a circle / sphere as a geometric
image. This is the invariant of the existence of difference (``there is
another quantity''), and corresponds to the real part g of the complex
inner product introduced later. (Which one is taken as the real part is,
by the principle of anonymity, also arbitrary.)

\textbf{Lemma (closure is not an axiom but a consequence of anonymity):}
{[}Interior{]} An end (boundary) is a privileged point with a neighbor
only on one side, and violates anonymity. Hence a system preserving
anonymity cannot have ends and closes --- moreover, it closes
periodically (rotationally symmetric). Therefore ``closed system'' is
not an independent axiom but is derived as a consequence of anonymity
(Axioms 1, 2). The closure assumption placed as an independent axiom in
the old version is removed in this revision.

\textbf{Note (mutual underivability of Axioms 4 and 5):} {[}Interior{]}
The product (Axiom 4, ω) and the sum of squares (Axiom 5, g) cannot be
derived from each other. This is not a technical matter: ``whether it
is'' and ``whether it differs'' are two root distinctions that stand
independently under anonymity. If one followed from the other, the limit
of anonymity would be wider than two.

\begin{center}\rule{0.5\linewidth}{0.5pt}\end{center}

\subsection{§3 Derivation of discreteness (discreteness is not an axiom
but a derived
quantity)}\label{derivation-of-discreteness-discreteness-is-not-an-axiom-but-a-derived-quantity}

The discreteness of the system is not posited as an axiom. It is derived
in two stages.

\textbf{Stage one (fundamental discreteness: from Axiom 2).}
{[}Interior{]} In the interior, 1 does not stand alone (Axiom 2). That
is, the unit ``1'' for subdividing a continuum, lacking contrast, cannot
be distinguished from 0. A continuum (continuous cardinality) stands
only when ``a unit that stands alone can be subdivided infinitely,'' but
that unit does not stand. Hence what can be constructed in the interior
is only the repetition of binary presence/absence judgments = doubling /
halving (×2, ÷2), and counting is discrete (integer). The base being 2,
integrality, and the half-bit introduced later all descend from this ``1
does not stand alone.''

\textbf{Stage two (discreteness of the mode number: from boundedness).}
{[}Discrete structure{]} By the lemma of §2, the boundedness of Axiom 5
appears as a \textbf{periodically closed system} (a system with no ends
that returns upon going around once). That it is not a string {[}0,L{]}
with ends is consistent with the lemma (ends violate anonymity). For a
standing wave to stand on a periodically closed system (circumference
L), the periodic boundary condition must be satisfied, and the
admissible wavelengths are quantized to integer fractions of the
circumference (λ\_n=L/n, n=1,2,3,\ldots). The spectrum is numbered by
the integer n.~The genesis of this integer mode number n is counting =
the genesis of dimension, and is the mechanism that breaks through the
ceiling (2) of Axiom 2 toward 3-or-more.

That is, discreteness is the name for what Axiom 5 does to Axioms 1--4,
and is not an independent principle.

\begin{center}\rule{0.5\linewidth}{0.5pt}\end{center}

\subsection{§4 Introduction of the fundamental quantity
N}\label{introduction-of-the-fundamental-quantity-n}

{[}Conserved quantity{]} The conserved quantity is not the linear
observables ν, λ but the squared quantity Σν²=R² (Axiom 5). The linear
quantities ν, λ are observables that swing reciprocally under the
conjugate relation νλ=k, and behind their product an invariant squared
quantity is conserved.

We introduce the fundamental quantity as the logarithm of a squared
quantity. The base is taken to be 2 by Axiom 2 (the interior binarity).
Attaching a factor 1/2 gives an integer sequence:

\[N \equiv \tfrac{1}{2}\log_2(\nu^2) = \log_2\nu .\]

{[}Fixed point{]} When ν=2\textsuperscript{n}, N=n (integer). N=0 corresponds to ν=1 (the
self-dual fixed point). {[}Interior{]} In the interior, ν=1 cannot be
distinguished from ν=0 (Axiom 2), but this does not negate 1 as a
fixed-point / integer-structure entity. N represents the whole integer
set $\mathbb{Z}$, symmetric in sign about 0. The sign of N corresponds to which
side of the conjugate pair one views from.

{[}Interior / Exterior{]} The imaginary unit i is merely a symbol
folding two real components with a 90° phase difference into complex
notation; the axioms of this system close over the reals (complex
numbers are not introduced). As noted later (§7 note), the complex
structure arises as a consequence of this real two-ness.

\begin{center}\rule{0.5\linewidth}{0.5pt}\end{center}

\subsection{§5 The additive identity of the conjugate relation (the g
side)}\label{the-additive-identity-of-the-conjugate-relation-the-g-side}

{[}Exterior (observer){]} Taking the base-2 logarithm of the conjugate
relation νλ=k turns the product into a sum:

\[\nu\lambda=k \;\xrightarrow{\ \log_2\ }\; N_\nu + N_\lambda = K \quad(K\equiv\log_2 k).\]

This is an identity, and the conserved quantity is K (the sum of the N
of the conjugate pair). Whichever of the three terms N\_ν, N\_λ, K is
chosen on the right-hand side,

\[N_\lambda = K - N_\nu,\qquad N_\nu = K - N_\lambda,\qquad N_\nu + N_\lambda = K\]

are preserved as identities, and the only thing produced by the
interchange is a relabeling --- a subtraction of logarithms (i.e.,
division). {[}Interior{]} Hence, without privileging any quantity,
anonymity is maintained.

{[}Exterior (observer){]} A change of scale (a change of k) corresponds
to a translation of the N axis, and scale invariance appears as
translational symmetry. At k=1, K=0 and N\_λ=−N\_ν (sign-reversal
symmetry about 0).

\textbf{Note (the identity generates no structure):} This additive
identity is the very definition of the logarithm, and by itself
generates no new structure. What generates structure is the fluctuation
Δ (the ω side) of the next section. This identity is merely a rephrasing
of the g side (distance, sum of squares).

\begin{center}\rule{0.5\linewidth}{0.5pt}\end{center}

\subsection{§6 The fluctuation Δ and uncertainty (the ω side, the side
that generates
structure)}\label{the-fluctuation-ux3b4-and-uncertainty-the-ux3c9-side-the-side-that-generates-structure}

What generates structure is the fluctuation Δ. To place this precisely,
we first kill two misreadings.

\subsubsection{6.1 Denial of hidden variables / observational
error}\label{denial-of-hidden-variables-observational-error}

{[}Discrete structure{]} ν is discrete (integer-valued) and has no
continuous fractional part. It is not that ``a continuous value ν=2.0±δ
really exists and appears integral due to the limits of observational
precision''; rather, the system is discrete, so ν is integral. This is
not a claim about reality but a consequence of discrete structure (see
the note in §0.5).

ν=2 means that ν fluctuates while remaining an integer, with its central
value (median) being 2. For example, the value can take the neighboring
integers 1, 2, 3, with center 2. The value itself jumps among integers.
This is a fluctuation intrinsic to the system, not an observational
error riding on a fixed true value.

{[}Interior{]} The interior quantity itself is an integer. Hence it
cannot take a fractional value such as 0.5. If fluctuation appears in
the interior, it appears only as an integer-step shift, i.e., ±1.

\subsubsection{6.2 ±1/2 is a posit, not a
derivation}\label{is-a-posit-not-a-derivation}

\textbf{{[}Negation{]}} One is tempted here to say ``the standard
deviation is ±1/2.'' But this report does not derive it. There is, in
this system, no basis on which to even construct the concept of a
standard deviation.

\textbf{{[}Reason{]}} Defining a standard deviation requires an
operation of observing the same object repeatedly and aggregating its
scatter. But ---

\begin{itemize}
\tightlist
\item
  {[}Interior{]} By anonymity, observations cannot be individuated from
  one another (this is the first, this is the second).
\item
  Since no time evolution is defined, there is no order in which to
  arrange observations.
\item
  It is not assumed that the true value is preserved the same between
  observations.
\end{itemize}

Hence the very phrase ``observe multiple times'' is self-contradictory
in this system. That single phrase smuggles in many hidden premises ---
individuation, temporal order, persistence of value.

\textbf{{[}Posit{]}} Therefore ±1/2 is not a computed value but a posit
(an assumption that it is presumably so). Whatever method one uses to
derive 1/2, one necessarily smuggles in one of the hidden premises
(individuation, temporal order, persistence). It is not that the
standard deviation cannot be computed exactly; it is that the very basis
for constructing the concept of a standard deviation is absent. This
report places ±1/2 as an assumption, not a derivation.

\subsubsection{6.3 Conservation law of the fluctuation (sum-zero is the
definition, constant-product is
derived)}\label{conservation-law-of-the-fluctuation-sum-zero-is-the-definition-constant-product-is-derived}

\textbf{Note (the fluctuations of §6.1 and §6.3 are different things):}
The ``the interior integer quantity fluctuates in integer steps (±1)''
stated in §6.1 and the half-bit fluctuation Δ=±1/2 placed in N-space in
this §6.3 are \textbf{separate posits; one is not derived from the
other}. No relation (identity, derivation, or approximation) is claimed
between them. §6.1 is a statement about the fluctuation of the interior
integer quantity; §6.3 is a statement about the fluctuation posited on
the side of the fundamental quantity N. This report treats them as
independent assumptions, and leaves their reconciliation or unification
as future work. Joining what is not joined inevitably smuggles in a
hidden premise (the same discipline as ±1/2 in §6.2), so this report
does not join them.

{[}Interior{]} We place the fluctuation on the side of the fundamental
quantity N (positing the fluctuation in N-space makes it scale-dependent
in the observable space, which is natural as a fluctuation of integer
counting). Each fundamental quantity is given a fluctuation Δ=±1/2 (a
half-bit; a posit; scale-invariant). The conjugate pair satisfies the
conservation of the sum N\_a+N\_b=K, and the fluctuations satisfy

\[\boxed{\ \Delta_a + \Delta_b = 0\ }\]

(which is taken as the definition).

{[}Exterior (observer){]} By the inverse map a=2\^{}\{N\_a+Δ\_a\},

\[k = ab = 2^{(N_a+N_b)+(\Delta_a+\Delta_b)} = 2^{K+0} = 2^K \quad(\text{constant})\]

is obtained, and the product of the observables ab is held constant
under fluctuations (the lower bound of minimal uncertainty). That is,
the uncertainty relation in the observable space (constant product) is
derived from the sum-zero of the fluctuations in N-space. Note that what
is preserved is the product of the observables ab, not the product of
the fluctuations δ\textsubscript{1}δ\textsubscript{2} (the latter depends on the allocation and is not
constant; see the appendix ``From the Sum-Zero of Fluctuations to
Uncertainty''). The minimal unit of uncertainty, 1/2, corresponds to the
posited half-bit ±1/2. Note that a single fluctuation ±1/2 and the
product of two fluctuations δ\textsubscript{1}δ\textsubscript{2}=1/2 (Axiom 4) are not the same
quantity; they are placed in correspondence as quantities belonging to
the same ω (area) side.

\begin{center}\rule{0.5\linewidth}{0.5pt}\end{center}

\subsection{§7 Main result: the logarithmic map of scale and the g/ω
separation (a limited
regime)}\label{main-result-the-logarithmic-map-of-scale-and-the-gux3c9-separation-a-limited-regime}

\textbf{{[}Note: on ``scale'' here.{]}} The object of flattening in this
section is R in Σν²=R², i.e., the scale (magnitude) of the system. This
is neither a curvature tensor nor the radius of a sphere. Σν²=R² is a
conservation law stating that the sum of squares equals R², and does not
require a circle or sphere geometry. This report does not claim
``flattening of curvature'' but states ``linearization (flattening) of
scale by the logarithm.'' Extension to a curvature tensor lies outside
the scope of this report (§8, §9). This report does not claim that the
continuum limit of this system extends to a curvature tensor; rather, it
regards it as likely to hold only in the discrete system.

\subsubsection{7.1 Complementarity of linearization in the two
spaces}\label{complementarity-of-linearization-in-the-two-spaces}

{[}Exterior (observer){]} The conserved quantity of Axiom 5 (an interior
invariant) appears to the observer as the sum of squares Σν²=R² (g,
scale R), and this is the conservation law representing scale. The
conjugate relation νλ=k is a product (nonlinear). In the N-space via the
logarithm, this relation is reversed. νλ=k is additive (N\_ν+N\_λ=K,
§5), i.e., a flat hyperplane, whereas the sum of squares Σν² is in
general a sum of exponentials Σ2\^{}\{2N\_n\} and is not linearized in
N.

\subsubsection{7.2 Flattening in the high-symmetry, anonymity
limit}\label{flattening-in-the-high-symmetry-anonymity-limit}

{[}Interior{]} Consider the fully symmetric limit (every axis can come
to the right-hand side, i.e., all ν\_n² are equal). This limit reduces
to all ν\_n² taking the same value v (the number of axes effectively
contracts to 2; this is also a reprise of the binarity of Axiom 2). Then
Σν²=(number of axes)·v, and the logarithm can act on the sum:

\[\log_2\!\left(\sum\nu^2\right) = \log_2(\text{number of axes}) + \log_2 v.\]

That is, the sum of squares representing scale is mapped onto a flat
linear additive relation:

\[\sum\nu^2 = R^2 \;\xrightarrow{\ \log_2,\ \text{full symmetry}\ }\; N + \tfrac12\log_2(\text{number of axes}) = N_R \quad(\text{an identity; flat}).\]

The source of scale is the square (power), and the logarithm linearizes
the power. Furthermore, since this system is scale-invariant, scale
dependence is flattened as a straight line along the translation. In
this limit, the non-uniform terms corresponding to the derivative terms
left by the Weyl transformation (§9) vanish under the assumption of full
symmetry (spatial uniformity).

However, the fully symmetric limit, as the price of making the
flattening hold exactly, almost exhausts the degrees of freedom on the g
side. That is, the substance of the flattening of §7.2 in itself is
thin. The substance of this report's claim lies not in the flattening
itself but in the g/ω separation of the next section §7.3 --- that even
when g degenerates to a single value, ω (quantization) remains.

\subsubsection{7.3 The heart of the map: separation of scale (g) and
quantization
(ω)}\label{the-heart-of-the-map-separation-of-scale-g-and-quantization-ux3c9}

This is the heart of the report. The flattening above does not ``throw
away the scale information.''

{[}Exterior (observer){]} The conserved quantity of Axiom 5 appears to
the observer as the sum of squares Σν², i.e., distance / metric g (the
real part of the complex inner product introduced later), the side of
the Born norm / normalization. The logarithmic map acts on this g side
and, in the fully symmetric limit, folds it into a flat coordinate. The
g side carries scale but not quantization. Scale resides wholly in g and
is flattenable (removable as a gauge).

{[}Interior{]} On the other hand, the fluctuation Δ = δ\textsubscript{1}δ\textsubscript{2}=1/2 (Axiom 4)
is area / the symplectic form ω (the imaginary part of the complex inner
product), the side that carries quantization. This ω is not crushed by
the logarithmic map and remains, on the flat N-surface, as the invariant
half-bit ±1/2. Even when g (the mean / scale) degenerates to a single
value v in the fully symmetric limit, ω (the fluctuation Δ) does not
degenerate.

That is, the logarithmic map is a \textbf{map that separates scale (g,
gauge) from quantization (ω, area, invariant)}. Here ``separation''
means, operationally, that the map flattens g (scale, sum of squares)
while leaving ω (the product δ\textsubscript{1}δ\textsubscript{2}) invariant. This is not a claim that
it is a canonical transformation strictly preserving the symplectic form
(that rigorization is future work). It pushes scale into the flat
surface and gauges it away, while purifying quantization as an invariant
on that flat surface. This is the substance of ``even when flattened in
the fully symmetric limit, the physics does not vanish'' --- the physics
(quantization) was, from the start, not in g (the mean) but resided in ω
(the fluctuation Δ).

\subsubsection{7.4 Structure of the map}\label{structure-of-the-map}

{[}Exterior (observer){]} In the self-dual fixed point k=1 and the fully
symmetric, anonymity limit:

\begin{itemize}
\tightlist
\item
  the scale-bearing system in the observable space (ν-space) (sum of
  squares Σν², scale R, the g side)
\item
  is mapped, via the logarithmic connection, onto a flat hyperplane in
  N-space (the affine structure of the linear sum ΣN=K),
\item
  and is exactly restored to the observable space by the inverse map
  ν=2\^{}N.
\item
  And the fluctuation Δ (the ω side) is preserved invariant before and
  after the map.
\end{itemize}

\textbf{Note (the complex structure as a consequence of real two-ness):}
{[}Interior / Exterior{]} Reinterpreting ν\textsubscript{1}, ν\textsubscript{2} as a radius ν and a
phase θ gives ν\textsubscript{1}=ν cosθ, ν\textsubscript{2}=ν sinθ, and the two are geometrically
isomorphic to the polar representation of z=ν e\^{}\{iθ\}. The sum of
squares ν²=ν\textsubscript{1}²+ν\textsubscript{2}²=\textbar z\textbar² corresponds to the real part g,
the area of the fluctuation δ\textsubscript{1}δ\textsubscript{2}=1/2 to the imaginary part ω, and the
cos↔sin rotation to the complex structure J. Complex numbers are not
placed in the axioms; the complex structure arises as a consequence of
real two-ness (the ``2'' of Axiom 2) and the sum of squares (Axiom 5).

\begin{center}\rule{0.5\linewidth}{0.5pt}\end{center}

\subsection{§8 Limitations}\label{limitations}

The scope of this result is strictly limited.

\begin{enumerate}
\def\labelenumi{\arabic{enumi}.}
\item
  \textbf{The result is a special solution, not a general one.} It holds
  only at k=1 (the self-dual fixed point) and in the fully symmetric,
  anonymity limit. Leaving this limit (k≠1, non-uniform, asymmetric),
  the non-uniform terms that vanished in §7.2 (corresponding to the
  derivative terms of the Weyl transformation) generally revive, and the
  flattening does not hold.
\item
  \textbf{±1/2 is a posit and has not been derived.} As in §6.2, this
  system has no basis on which to construct the concept of a standard
  deviation. Every computation deriving ±1/2 smuggles in a hidden
  premise. Note that the main result of this report (the g/ω separation
  of §7.3) presupposes this posit ±1/2. If the posit fails, the
  conclusion changes.
\item
  \textbf{No exhaustive search for counterexamples has been carried
  out.} It has not been shown that the result has no counterexamples.
\item
  \textbf{The relation to known frameworks is undetermined.} Whether the
  result remains a special case of known frameworks (Weyl
  transformation, conformal flatness, conformal cosmology, etc.) or
  exceeds them cannot be determined at present. In particular, how the
  g/ω separation of §7.3 differs from or agrees with known Kähler
  geometry / conformal flatness has not been rigorously compared.
\item
  \textbf{This report treats only the scale R, not a curvature tensor.}
  The R of Σν²=R² is the scale (magnitude) of the system, not a
  curvature tensor or the radius of a sphere. A rigorous proof including
  the differential structure (the curvature tensor) is not given, and
  extension to curvature / general curved spacetimes lies outside the
  scope of this report.
\item
  \textbf{No physical identification is made.} This system is an
  algebraic / geometric (informational) structure, and no physical
  interpretation (time, space, energy, particles, etc.) is attached to
  the observables / fundamental quantities. Isomorphism with a physical
  theory is a separate task.
\item
  \textbf{No relation or unification of the two fluctuations is given.}
  The integer-step fluctuation (±1) of §6.1 and the half-bit fluctuation
  Δ=±1/2 of §6.3 are separate posits, and no relation (identity,
  derivation, approximation) is claimed between them. The derivation of
  placing the fluctuation in N-space from the interior integer
  fluctuation is not carried out in this report. Unification is left as
  future work.
\end{enumerate}

\begin{center}\rule{0.5\linewidth}{0.5pt}\end{center}

\subsection{§9 Positioning relative to known frameworks (for
reference)}\label{positioning-relative-to-known-frameworks-for-reference}

Computing physical quantities in a curved spacetime is generally
difficult. One existing approach is the Weyl transformation
\(g_{\mu\nu}\to\Omega^2 g_{\mu\nu}\); but the standard action is not
invariant under Weyl rescaling, and the curvature scalar transforms as

\[R \to e^{-2\sigma}\left(R - 6\,\Box\sigma - 6\,g^{\mu\nu}\partial_\mu\sigma\,\partial_\nu\sigma\right),\]

so the curvature information remains as derivative terms of the
rescaling function σ. That is, even after flattening, the curvature does
not entirely vanish. Conformal flatness always holds locally in two
dimensions but not for general metrics in higher dimensions.

From a starting point different from these known frameworks --- only a
few axioms --- this report presents the structure that, in a limited
regime, a quantity representing scale (not a curvature tensor) is mapped
onto a flat structure leaving no derivative terms, and quantization (ω)
is preserved invariant under the map. What this report treats differs in
scope from the Weyl transformation, which rescales a curvature tensor:
it is a logarithmization of scale. It is neither a rehash of Weyl nor a
claim to treat curvature. Extension to curvature is left as future work.
No generalization is claimed. The possibility that it is a special case
of a known framework is not excluded either.

\begin{center}\rule{0.5\linewidth}{0.5pt}\end{center}

\subsection{Supplements}\label{supplements}

This report has appendices that concretize points compressed in the main
text and the ``arising structure.'' All are \texttt{isSupplementTo} this
report.

\begin{enumerate}
\def\labelenumi{\arabic{enumi}.}
\tightlist
\item
  \textbf{A Magic-Coin Toy Model} --- visualizes the limit of anonymity
  (0, 1, 2) and the denial of hidden variables of §0.5 as a tabletop
  magic. Concept DOI:
  \href{https://doi.org/10.5281/zenodo.20741264}{10.5281/zenodo.20741264}
\item
  \textbf{Multidimensional Structure from the Logarithmic Fundamental
  Quantity} --- §3, §5, §7. Exhibits the (n−1)-dimensional hyperplane
  defined by the additive identity ΣN\textsubscript{i}=K (the minimal example of the
  ``arising structure''). Concept DOI:
  \href{https://doi.org/10.5281/zenodo.20741712}{10.5281/zenodo.20741712}
\item
  \textbf{From the Sum-Zero of Fluctuations to Uncertainty} --- §6.3.
  Makes explicit the direction in which the constant product of
  observables is derived from the sum-zero in N-space (what is preserved
  is the product of observables, not the product of fluctuations).
  Concept DOI:
  \href{https://doi.org/10.5281/zenodo.20742277}{10.5281/zenodo.20742277}
\end{enumerate}

\begin{center}\rule{0.5\linewidth}{0.5pt}\end{center}

\subsection{§10 Conclusion}\label{conclusion}

Starting from only a few axioms (difference; the limit of anonymity = 2;
fluctuation; conservation of the product; conservation of the sum of
squares), strictly distinguishing the system-interior (the unknowable)
from the system-exterior (observation), and placing the objective
anchors not on reality but on fixed points / conserved quantities, this
report records, as a limited structure, that when the fundamental
quantity is introduced as the logarithm of a squared quantity N=log\textsubscript{2}ν,
in the extremely high-symmetry, anonymity limit, the sum of squares
representing scale (Σν²=R², g) can be mapped via the logarithmic
connection onto a flat hyperplane (affine structure), and that this map
separates scale (g) from quantization (ω, the fluctuation Δ), preserving
quantization as an invariant on the flat surface. Here R is the scale of
the system, not a curvature tensor.

This is not claimed as the discovery of a general fact; it is recorded
as a result within a limited scope, with the posited nature of ±1/2 and
the incompleteness of the counterexample search explicitly stated.
Whether this structure is a special case of a known framework, and
whether a generalization is possible, are left to future study.

\begin{center}\rule{0.5\linewidth}{0.5pt}\end{center}

\emph{This report is a preliminary, limited record, and contains no
definitive claim regarding the generality, rigor, or novelty of the
result.}
\end{document}
